% Publication gate: retain the 87-matrix Lean claim only after the root
% worker obtains the complete PASS receipt for all 87 FourMatrix239 modules.
\section{Why four missing counts can also be recovered}
The neighbour counts at any 508 of the 512 vertices of $Q_9$ determine
a binary vertex selection uniquely. If the selection lies on a known
parity side, any 252 of the 256 opposite-side counts suffice. The counts
on the selected side are automatically zero.

To compare two selections, subtract their indicators. Their difference
$f$ takes values $-1,0,1$. If the observed counts agree, the neighbour
sums of $f$ can be nonzero only at the four missing positions. When
those positions have mixed parity, splitting $f$ by parity reduces
each part to at most three missing counts, already covered by the
preceding result.

For four positions of one parity, apply the inverse formula at their
antipodes. If $w$ lists the four missing count differences, the resulting
integer equations are
\[
Mw=315s,\qquad s\in\{-1,0,1\}^4.
\]
The diagonal entries of $M$ are 128. Its other entries are $-16,8,-10,35$
when the corresponding missing positions are at distance $2,4,6,8$.

Only 35 is odd. Join positions at distance eight. Two such neighbours
of one position have distance two from each other, so this graph has
no triangle. There are seven possible graph types on four vertices.
Parity, subtraction of equations and divisibility settle five types
directly. The other two reduce to seven triples of coefficients and
five small block matrices. In each case, elimination produces
$Dx=315N$; the integer $D/\gcd(D,315)$ must divide $N$, while the allowed
right-hand sides make $|N|$ too small unless $N=0$. The full paper gives
the matrices and numerical bounds. These arguments force $w=0$, and
the inverse then forces the entire selection difference $f$ to vanish.

A further result permits up to eight missing counts when no pair of
missing positions has Hamming distance eight. In that case every
off-diagonal coefficient in the antipodal equations is even, with
absolute value at most 16. Parity forces every right-hand side to zero.
If $k\le8$ counts are missing and $M$ is the largest absolute missing
count difference, a row attaining it gives
\[
128M\le16(k-1)M\le112M,
\]
so $M=0$. This is a written deduction from the inverse. The condition
on the missing positions is required by this argument.

Eight is sharp under that condition. The nine neighbours of one vertex
have pairwise distance two, yet losing all nine counts conceals whether
that vertex is selected. The unrestricted result above remains four
missing counts; whether every set of five through eight missing counts
can be tolerated is unresolved. Any failure in that range must contain
a pair of missing positions at distance eight.

The classification was also checked independently. Grouping the nine
coordinate columns by their three-bit patterns gives 1149 valid
ordered descriptions, 105 cube-isometry classes of four-point sets,
and 87 distinct distance matrices. A second enumeration checks all
2,731,135 triples after fixing one vertex at zero and obtains exactly
the same 87 matrices. All 7047 matrix/right-hand-side combinations have
no nonzero integer solution~\cite{research239}.

All 87 matrix implications passed Lean checking, including their
adjugate identities and divisibility calculations. The connection from
arbitrary cube positions to this catalogue, and the parity reduction,
remain written arguments with independent computational support. The
complete four-count recovery statement has not been assembled as one
Lean theorem.

Thus two different selections must alter at least five neighbour
counts. Changing one selected vertex alters nine, so the smallest
possible number remains between five and nine. Recovery alone does
not ensure that a selection is a connected induced path. No longer
snake or new global bound follows from this result.
